% Created 2026-06-05 Fri 17:33
% Intended LaTeX compiler: pdflatex
\documentclass[presentation]{beamer}
\usepackage[utf8]{inputenc}
\usepackage[T1]{fontenc}
\usepackage{graphicx}
\usepackage{longtable}
\usepackage{wrapfig}
\usepackage{rotating}
\usepackage[normalem]{ulem}
\usepackage{amsmath}
\usepackage{amssymb}
\usepackage{capt-of}
\usepackage{hyperref}
\makeatletter\def\input@path{{./}{./shared/}{./shared/themes/}}\makeatother
\usepackage{beamerthemequeer}
\usepackage{shared-preamble}
\usetheme{default}
\author{Matt Alexander}
\date{June 6, 2026}
\title{A (Pr)operad of (Pr)operad Algebras}
\hypersetup{
 pdfauthor={Matt Alexander},
 pdftitle={A (Pr)operad of (Pr)operad Algebras},
 pdfkeywords={},
 pdfsubject={},
 pdfcreator={Emacs 30.1 (Org mode 9.7.11)}, 
 pdflang={English}}
\begin{document}

\maketitle
\begin{frame}[label={sec:org2bbcbda}]{Main Results}
\begin{proposition}{The Operad of Operad Algebras}{operad-of-operad-alg-gen}
Let $\mathcal{V}$ be a symmetric monoidal category with all small coproducts such that the monoidal product distributes over all \emph{finite} coproducts, and let $\mathcal{O}$ be a coloured operad in $\mathcal{V}$. Then for any small category, $\mathcal{C}$, there exists an operad $\mathcal{O}^{\mathcal{C}}$ in $\mathcal{V}$ whose algebras are the functors 
\begin{equation*}
F: \mathcal{C} \to \mathcal{O} \text{-Alg}.
\end{equation*}
\end{proposition}
\end{frame}
\begin{frame}[label={sec:org1908d41}]{Main Results (Properad Version)}
\begin{proposition}{Properad of Properad Algebras}{prop-of-prop-algs}
Let $\mathcal{P}$ be a coloured properad in a closed monoidal category with all small coproducts $\mathcal{V}$, where the monodial product distributes over all \emph{finite} coproducts. For any small category $\mathcal{C}$, there exists a properad $\mathcal{P}^{\mathcal{C}}$ whose algebras are the functors
\begin{equation*}
F: \mathcal{C} \to \mathcal{P}\text{-Alg}.
\end{equation*}
\end{proposition}
\end{frame}
\begin{frame}[label={sec:orgbd9da24},fragile]{Generators and Relations}
 \newthought{Between objects in $\mathcal{C}$: }
\[
\mathcal{O}^{\mathcal{C}}_{\text{gen}}[A_{\alpha} \mid B_{\alpha}] = \bigsqcup_{f: A \to B} I
\]
when \(A \ne B\).

\newthought{Copies of $\mathcal{O}^{\mathcal{C}}$: }
\[
\mathcal{O}^{\mathcal{C}}_{\text{gen}}[A_{\alpha_{1}} \ldots A_{\alpha_{n}} \mid A_{\beta}] = \mathcal{O}[\alpha_{1} \ldots \alpha_{n}]
\]
\newthought{Endomorphisms: } 
\[
\mathcal{O}^{\mathcal{C}}_{\text{gen}}[A_{\alpha} \mid A_{\alpha}] = \mathcal{O}[\alpha \mid \alpha] \sqcup \bigsqcup_{f: A \to A} I
\]
\end{frame}
\begin{frame}[label={sec:org4c6acf1}]{Explicit Form (infinitary distributivity)}
In the case that \(\otimes\) distributes over \emph{all small coproducts}
\[
\mathcal{O}^{\mathcal{C}}[\bar{A}_{\bar{\alpha}} \mid B_{\beta}] = \coprod_{f^i: A^i \to B} \mathcal{O}[\bar{\alpha} \mid \beta]
\]
\bigskip

Similarly, there's a more involved properadic version
\[
\mathcal{P}^{\mathcal{C}}[\bar{A}_{\bar{\alpha}} \mid \bar{B}_{\bar{\beta}}] = \coprod_{\text{Alternating chains}} \mathcal{P}[\bar{\alpha} \mid \bar{\beta}]
\]

\[
A \to B \rightrightarrows C_{i} \to D \rightrightarrows \ldots \to Z
\]
\end{frame}
\begin{frame}[label={sec:orgcb80933}]{Precursors}
\begin{block}{\emph{Coloured Operads}, Yau (2017)}
\begin{itemize}
\item The construction of \(\mathcal{O}^{\mathcal{C}}\) is left as an exercise.
\end{itemize}

\bigskip
\end{block}
\begin{block}{\emph{Operads for algebraic quantum field theory}, Benini, Schenkel, Woike (2021)}
\begin{itemize}
\item \(\mathcal{O}^{\mathcal{C}}\) is constructed in the case that \(\mathcal{O}\) is the associative operad:
\end{itemize}
\[
\operatorname{Assoc}^{\mathcal{C}}
\]
\end{block}
\end{frame}
\begin{frame}[label={sec:org6cfde49},fragile]{Orthogonal Categories}
 \begin{block}{Definition}
Let \(\mathcal{C}\) be a small category. An \emph{orthogonality relation} \(\perp\) on \(\mathcal{C}\) is a relation on pairs of morphisms in \(\mathcal{C}\) such that if \(f \perp g\)
then \(f\) and \(g\) have the \emph{same codomain}, and are required to satisfy

\begin{itemize}
\item If \(f \perp g\) then \(g \perp f\).

\item If \(g \perp h\) then \(fg \perp fh\) for all composable \(f\).

\item If \(g_1 \perp g_2\) then \(g_1 h \perp g_2 k\) for all composable \(h,k\).
\end{itemize}

\bigskip 
An \emph{orthogonal category} is a small category \(\mathcal{C}\) with an orthogonality relation \(\perp\).

\newthought{Intuition: } \(f \perp g\) means that the regions corresponding to \(f\) and \(g\) are \emph{causally disjoint}
\end{block}
\end{frame}
\begin{frame}[label={sec:org9eb05e9}]{Causality}
In order to be a physical theory, an AQFT must satisfy the \emph{causality condition}:

\begin{center}
\begin{tikzcd}[ampersand replacement=\&]
F(A) \otimes F(B) \arrow[r, "F(f) \otimes F(g)"] \arrow[d, swap, "F(f) \otimes F(g)"] \& F(C) \otimes F(C) \arrow[r, "\tau"] \& F(C) \otimes F(C) \arrow[d, "\mu"] \\
F(C) \otimes F(C) \arrow[rr, swap, "\mu"]  \& \& F(C)
\end{tikzcd}
\end{center}

\bigskip

whenever \(f \perp g\).
\end{frame}
\begin{frame}[label={sec:orge8eb1b8}]{Introducing Homotopy}
To handle gauge theories we need homotopy:
\begin{itemize}
\item BV-BRST formalism
\item Aharonov-Bohm effect (Wilson loops are non-local)
\end{itemize}

\bigskip

We can't expect causality to hold strictly

\bigskip
\begin{block}{Homotopy Quantum Field Theory, Yau (2019)}
\begin{itemize}
\item Wraps the entire QFT in homotopy:
\end{itemize}
\[
W(\operatorname{Assoc}^{\mathcal{C}})
\]
\end{block}
\end{frame}
\begin{frame}[label={sec:org8f89d3b}]{Pointwise A\textsubscript{\(\infty\)}-algebras}
\begin{block}{Homotopy QFT}
Picking out a monoid at any point in the category \(\mathcal{C}\)
\[
\operatorname{Assoc} \to \mathcal{O}^{\mathcal{C}}
\]
induces an adjunction
\[
\operatorname{Alg}(\operatorname{Assoc}_{\infty}) \rightleftarrows \operatorname{Alg}(W(\mathcal{\operatorname{Assoc}}^{\mathcal{C}})),
\]

\bigskip 

which picks out the corresponding A\textsubscript{\(\infty\)}-algebra at that object.
\end{block}
\end{frame}
\begin{frame}[label={sec:orgad321a8}]{Alternatives}
Using the \(\mathcal{O}^{\mathcal{C}}\) construction, we can create a model for AQFTs of the form
\[
\operatorname{Assoc}_{\infty}^{{\mathcal{C}}}.
\]
This gives us finer control over the homotopy \emph{in the algebras}.

\begin{itemize}
\item This is an easier model to perform computations in
\end{itemize}

\bigskip
\begin{block}{Properads}
Structures like Hopf algebras, L\textsubscript{\(\infty\)}-bialgebras, and quantum groups are \emph{properadic}.
\begin{itemize}
\item We can take our properadic construction and then inject homotopy through
\begin{itemize}
\item Vallette's Properadic Koszul Duality
\item Kontsevich's Graph Complexes
\end{itemize}
\end{itemize}
\end{block}
\end{frame}
\begin{frame}[label={sec:org8d529b6},fragile]{Hopf-Laplace Algebras}
 A \emph{bialgebra} is a unital associative \$k\$-algebra, \(B\), with maps 

\medskip

\begin{itemize}
\item \fbox{Comultiplication: } \(\Delta: B \to B \otimes B\)

\item \fbox{Counit: } \(\varepsilon: B \to k\)
\end{itemize}

\bigskip

satisfying the appropriate compatibility conditions between the algebraic and coalgebraic structure.

A \emph{Laplace pairing} on a bialgebra is a bilinear form \(\gen{\cdot, \cdot}: B \otimes B \to k\) satisfying

\[
\gen{x, yz} = \sum \gen{x_{(1)}, y} \gen{x_{(2)}, z}
\]
\[
\gen{xy, z} = \sum \gen{x, z_{(1)}} \gen{y, z_{(2)}}
\]
\[
\gen{1, x} = \gen{x, 1} = \varepsilon(1)
\]

\newthought{Intuition: } Laplace pairings behave like \emph{correlation functions} from quantum field theory!
\end{frame}
\begin{frame}[label={sec:org0274385}]{Properad of Laplace Hopf AQFTs}
Let \((\mathcal{C}, \perp)\) be an orthogonal category. A \emph{Laplace Hopf algebraic quantum field theory} is a functor
\begin{align*}
F: \mathcal{C} \to \mathcal{LB}(\mathcal{V})
\end{align*}
where \(\mathcal{LB}(\mathcal{V})\) is the category of Laplace bialgebras on a symmetric monoidal category \(\mathcal{V}\), such that the following \textbf{causality diagram} commutes:

\begin{center}
\begin{tikzcd}[column sep = 3cm, ampersand replacement=\&]
F(A)^{\otimes m} \otimes F(B)^{\otimes n} \ar[r, "F(f)^{\otimes m} \otimes F(g)^{\otimes n}"] \ar[d, swap, "\gen{-,-}_m \otimes \gen{-,-}_n"] \& F(C)^{\otimes m} \otimes F(C)^{\otimes n} \ar[d, "\gen{-,-}_{m + n}"] \\
k \otimes k \ar[r, swap, "\mu_k"]  \& k 
\end{tikzcd}
\end{center}
commutes whenever \(f: A \to C\) and \(g: B \to C\) satisfy \(f \perp g\), where \(\gen{-,-}_{m+n}\) is the (m+n)-arity correlation function on \(F(C)\).
\end{frame}
\end{document}
